
Large language models (LLMs) have recently gained significant attention in  natural language generation (NLG)~\citep{touvron2023llama,katz2023gpt,openai2023gpt4,chowdhery2022palm}.
Trained on vast amounts of data, they exhibit impressive abilities in generating human-like responses.
As such advances invariably lead to wider adoption of LLM for language generation tasks, such as question-answering (QA), it is crucial to quantify their uncertainty.

Uncertainty quantification (UQ) is well-studied within machine learning. 
A reliable measure of uncertainty is important to decide when to trust a model. 
When a model shows high uncertainty or returns low-confidence predictions\footnote{See \cref{sec:background:uvsc} for a discussion of uncertainty vs. confidence.}, the input should either be rejected or yielded to further evaluation~\citep{mcdp}.
This is an area of active research typically called selective classification~\citep{Chow1970OnTradeoff,lin2022scrib,Geifman2017SelectiveNetworks}. 
Similarly, in NLG, one might refuse the generation $\mathbf{s}$ by the LLM to the input $x$ with high uncertainty/low confidence. 
Selective generation basing on uncertainty estimates could potentially improve the decision making process, especially for high-stakes applications such as medical or legal question-answering.

While UQ has been an important topic for classical machine learning tasks like classification or regression~\citep{Lakshminarayanan2017SimpleEnsembles,mcdp,DBLP:conf/icml/Hernandez-Lobato15b,ABDAR2021243}, it bears specific challenges for NLG and has attracted limited attention until recently~\citep{Lin2022TeachingMT,kuhn2023semantic,malinin2021uncertainty}. 
Numerous challenges hinder the direct application of UQ methods from classification or regression. 
First, the output space has forbiddingly high dimensionality, rendering most measures, such as the exact entropy of predicted probabilities, unfeasible. 
Furthermore, distinct token sequences may convey identical meanings, necessitating an effective uncertainty measure to operate in the semantic meaning space~\citep{kuhn2023semantic}. 
Lastly, many existing LLMs are black-boxes served via APIs, implying that end-users typically lack white-box access. 
Even when such access is available, many LLMs are too large to run for most. 
Note that such considerations are orthogonal to the problem of overconfidence: We need a ranking-wise informative confidence/uncertainty measure before the problem of overconfidence even appears (see discussion in \cref{sec:related}).


In this study, we explore the problem of uncertainty quantification in NLG with black-box LLMs. 
Unlike some prior research~\citep{DBLP:journals/corr/abs-2012-14983,Lin2022TeachingMT,kadavath2022language,kuhn2023semantic}, we focus on the more realistic scenario where we only possess access to the generated text, rather than the numerical outputs such as token-level logits. 
We first introduce a set of measures designed to assess the uncertainty of the input and the confidence of each generation -  
The main idea entails estimating the uncertainty/confidence from multiple generations from the LLM. 
Then, we demonstrate the application of these uncertainty estimates in the context of free-form QA.
QA datasets are used for the simplicity of evaluation, as completely open-ended NLG tasks require expensive human evaluation.
The main contributions of this paper are summarized as follows:
\begin{itemize}
    \item We investigate uncertainty quantification for black-box LLMs, a previously under-explored topic, and assess its value in the downstream task of selective natural language generation. 
    \item We put forward several simple yet effective techniques for estimating uncertainty associated with input data and determining the confidence level of each individual generated response.
    \item Through extensive experiments on several popular LLMs and question-answering datasets, we observe that proposed measures demonstrate a high level of effectiveness in pinpointing challenging (uncertain) questions and predicting the quality of their corresponding answers.
\end{itemize}
